\section{Division coefficients and monomial ideals}\label{sec:reduction}

We first reduce the Hilbert-function problem to identities among commuting
linear forms. Their coefficients will lie in a division algebra, but the
linear forms themselves will commute. This distinction lets us use
commutative algebra to obtain a basis and division-algebra linear algebra
to count its least exponents.

Throughout, a field \(k\) extending the original \(\C\) is identified with
a central subfield of every \(k\)-algebra in use. In
\(\Delta[x_1,\ldots,x_n]\), the variables \(x_j\) are central and have
degree one, and elements of \(\Delta\) have degree zero.

\begin{lemma}[Ordered-basis criterion]\label{lem:ordered-basis}
Let \(f_1,\ldots,f_n\) be a homogeneous regular sequence in
\(S=\C[x_1,\ldots,x_n]\), with \(\deg f_i=a_i>0\).
Suppose that a field extension \(k/\C\) and a finite-dimensional central
division \(k\)-algebra \(\Delta\) have the following property.
There is a finite collection of pairwise commuting elements of
\(\Delta[x_1,\ldots,x_n]_1\), containing \(x_1,\ldots,x_n\), whose
generated commutative \(k\)-algebra \(A\) contains elements
\(t_1,\ldots,t_n\) of degree one satisfying
\begin{equation}\label{eq:triangular}
 t_i^{a_i}
 =c_i f_i+\sum_{l<i}B_{il}f_l+\sum_{d>i}C_{id}t_d,
 \qquad 1\le i\le n,
\end{equation}
where \(c_i\in k^\times\) and all terms are homogeneous, with
\(B_{il},C_{id}\in A\).
Then \(A\) is finite over \(k[x_1,\ldots,x_n]\), and
\[
 \{t_1^{\alpha_1}\cdots t_n^{\alpha_n}:\alpha\in\Z_{\ge0}^n\}
\]
is a basis of \(\Delta[x_1,\ldots,x_n]\) as a left
\(\Delta\)-vector space.
\end{lemma}

\begin{proof}
The regular sequence gives
\begin{equation}\label{eq:regular-hilbert}
 \Hilb_{S/(f_1,\ldots,f_n)}(t)
 =\frac{\prod_{i=1}^n(1-t^{a_i})}{(1-t)^n}
 =\prod_{i=1}^n(1+t+\cdots+t^{a_i-1}).
\end{equation}
Thus the \(f_i\) have no common projective zero, and this remains true
after any field extension.

Put \(N=\Delta[x_1,\ldots,x_n]\), considered as a right \(A\)-module.
The inclusions
\[
 k[x_1,\ldots,x_n]\subseteq A\subseteq N
\]
make \(A\) a submodule of the finite free \(k[x]\)-module \(N\).
Consequently \(A\) is finite over \(k[x]\), and both \(A\) and the finite
right \(A\)-module \(N\) are Noetherian.

We claim that \(t_1,\ldots,t_n\) is an \(N\)-regular sequence.
First, \(A/(t_1,\ldots,t_n)\) is zero dimensional. Indeed, at a geometric
point where all \(t_j\) vanish, the identities
\eqref{eq:triangular}, read in increasing order of \(i\), give
\(f_1=\cdots=f_n=0\). Equation~\eqref{eq:regular-hilbert} then forces
all \(x_j=0\), and finiteness over \(k[x]\) gives the asserted
zero-dimensional quotient.

Let \(\mathfrak m=A_{>0}\), the graded maximal ideal. The central
variables \(x_1,\ldots,x_n\) form an \(N\)-regular sequence because
\(N=\Delta[x]\) is free over \(k[x]\). Localization preserves the
successive injections. Their quotient is \(\Delta\), on which every
positive-degree element of \(A\) acts trivially, so it remains nonzero
on localization. Thus the sequence is also \(N_{\mathfrak m}\)-regular.
The support dimension of \(N_{\mathfrak m}\) is at most \(n\) by
finiteness over \(k[x]\). It therefore equals its depth \(n\).
The system-of-parameters theorem for finite Cohen--Macaulay modules
now implies that \(t_1,\ldots,t_n\) is regular on \(N_{\mathfrak m}\);
see \cite[Proposition~10.103.4, Tag 00N6]{Stacks}.
It is regular on the graded module \(N\) as well. To see the
localization step explicitly, a nonzero homogeneous element of any
graded kernel or quotient cannot be annihilated by an element of
\(A\setminus\mathfrak m\): the nonzero constant term of that element
would survive in the lowest-degree component.

Since all the \(t_j\) have degree one,
\[
 \Hilb_{N/\sum_j Nt_j}(t)
 =(1-t)^n\Hilb_N(t)=\dim_k\Delta.
\]
Its degree-zero part is already \(\Delta\). Hence, in every positive
degree, \(N_d=\sum_jN_{d-1}t_j\). Induction on degree, using
commutativity of the \(t_j\), shows that
\[
 \{\,t^\alpha=t_1^{\alpha_1}\cdots t_n^{\alpha_n}:
          \alpha\in\Z_{\ge0}^n\,\}
\]
spans \(N\) with coefficients on the left in \(\Delta\).
In degree \(d\) its cardinality is
\(\binom{n+d-1}{n-1}=\dim_\Delta N_d\); it is therefore a left
\(\Delta\)-basis.
\end{proof}

The full basis permits a single exponent comparison for every element of
an ideal. We now extract the Hilbert-function consequence; no further
geometric input is needed.

\begin{proposition}[Monomial extraction]\label{prop:monomial-reduction}
Let \(f_1,\ldots,f_n\), \(k\), \(\Delta\), \(A\), and
\(t_1,\ldots,t_n\) satisfy the hypotheses of
Lemma~\ref{lem:ordered-basis}. For every homogeneous ideal
\(I\subseteq S=\C[x_1,\ldots,x_n]\) containing the \(f_i\),
there is one monomial ideal \(J\subseteq S\) such that
\[
 x_i^{a_i}\in J\quad(1\le i\le n),\qquad
 \dim_\C(S/J)_d=\dim_\C(S/I)_d\quad(d\ge0).
\]
\end{proposition}

\begin{proof}
Put \(N=\Delta[x_1,\ldots,x_n]\), with the full left
\(\Delta\)-basis of Lemma~\ref{lem:ordered-basis}.

Extend \(I\) to the two-sided ideal \(\Delta I\) of \(N\). For each
nonzero homogeneous element of this ideal, take the least exponent
in its \(t\)-basis expansion under lexicographic comparison of the
coordinates in the order
\[
 \alpha_n,\alpha_{n-1},\ldots,\alpha_1.
\]
Let \(E\subseteq\Z_{\ge0}^n\) be the set of all such exponents.
Right multiplication by \(t_j\) translates every exponent by the
same vector \(\mathbf e_j\), without moving any left coefficient.
Thus \(E+\mathbf e_j\subseteq E\). In each degree, elimination over
the division algebra \(\Delta\) shows that the number of exponents
in \(E\) is the left dimension of \((\Delta I)_d\).

For each \(i\), Equation~\eqref{eq:triangular} puts
\[
 t_i^{a_i}-\sum_{d>i}C_{id}t_d
\]
in \(\Delta I\). Expand the \(C_{id}\) in the \(t\)-basis with
coefficients on the left. Every exponent in a term \(C_{id}t_d\)
has a positive coordinate of index greater than \(i\), and hence is
strictly larger than \(a_i\mathbf e_i\) in the chosen order.
It follows that \(a_i\mathbf e_i\in E\).

The upward-closed set \(E\) defines a single monomial ideal
\[
 J=\langle\,x^\alpha:\alpha\in E\,\rangle\subseteq S.
\]
It contains the required pure powers. Finally,
\[
 N/\Delta I\simeq \Delta\otimes_\C(S/I)
\]
as graded left \(\Delta\)-vector spaces. The complementary exponent
count in every degree gives the asserted Hilbert function.
\end{proof}

Our task is now to construct \eqref{eq:triangular}, starting with
row \(n\) and proceeding down to row \(1\). At a stage of degree
\(b=a_i\), the constraints on a new sum of \(b\)-th powers involve
the degree-\(b\) part of a formal ideal. We therefore encode earlier,
possibly higher-degree rows as consequences of linear and quadratic
relations, whose generators still have degree at most \(b\).
Section~\ref{sec:low-relations} defines these constraints and proves
the simple-factor property needed for their parameter space.
